\section{20 августа 1914 г.}

\lp{20260604_194334196}

Дорогой мой! Не имею больше писем от вас и очень беспокоюсь. Хочется знать как
вы себя чувствуете. С кем вы? Сколько времени думаете оставаться в Цюрихе? Я
все время мыслями с вами. Сегодня вечером мне мешали писать вам, а сейчас уже
первый час ночи.

\section{21 августа 1914 г.}

\lp{20260604_194334196}

Ужасное положение в том отношении, что даже поделиться с вами всем, что сейчас
волнует, нельзя. И когда все это кончится? И

\lp{20260604_194350913}

\noindent
когда мы с вами увидимся\textellipsis

А главное: я ничего не знаю про вас. Напишите о себе!

У нас сейчас живет Ирочка. Занимает она Колину комнату, а Коля пока в вашей
спальне. У Зяськи сейчас нет квартиры, и он с Борей живут у родителей. Шура в
Петербурге, а Виктория осталась у брата своего в деревне.

Родители ваши оба так же как и всегда бодры и здоровы. Сейчас мамаша очень
много горячится по поводу событий.

\lp{20260604_194350913}

Зяська благодушен и очень хорошо выглядит после Финляндии. Боря носит дедушкины
костюмы и похож на 25-летнего юношу. Верочка еще развилась, и в ее рассуждениях
появилась даже зрелость. Колины и Лёлины дела все также недостаточно выяснены.
Во всяком случае долгов они пока еще не выплачивают. Лёля все время имеет
работу, но весь доход уходит в расход

\lp{20260604_194334196}

\noindent
уравнивать приход с расходами.

Маня вернулась и принялась за старое: т.е. за туалеты -- свои и других дам. Папа
все еще не вернулся, и мы о нем очень беспокоимся. Сегодня мама получила через
какое-то консульство извещение, что он в Берлине и просит денег. Не напишите ли
вы ему туда два слова. Из Швейцарии, вероятно, письмо дойдет, а отсюда нет.

До свиданья, мой дорогой. Страшно не достает вас! Обнимаю. Ваша Анюта.

Это мое пятое письмо с момента получения телеграммы.

